\section{Related Work}

This work is inspired by the role that %
Torch \citep{collobert2002torch},
Theano \citep{bergstra2010theano,bergstra2011theano,al2016theano},
Chainer \citep{tokui2015chainer}, and others
played in the development in deep learning by providing powerful abstractions. A similar transformation is emerging with higher-level pipelines of LMs, and
we are seeking to offer a solid conceptual framework and programming abstractions for %
what we call \textit{foundation model programming}.
We draw on differentiable programming \citep{wang2018differentiable} but applied to LM calls rather than neural networks, and borrow syntactic elements from PyTorch \citep{pytorch2019}.

In-context learning (\citealt{mccann2018natural,radford2018improving,brown2020language}) is a key mechanism for foundation model programming. A growing body of work has revealed that, especially with instruction tuning \citep{ouyang2022training}, we can elicit sophisticated behavior via prompting
\citep{%
wei2022chain,
wang2022self,
press2022measuring,
yao2022react,
khot2022decomposed,
madaan2023self}.
Similarly, forms of weak supervision that would normally require
task-specific
\citep{%
khattab2021baleen,
khattab2021relevance}
or
hand-built 
\citep{%
ratner2016snorkel,hancock2018training}
heuristics are now done by LMs
\citep{
wang2022self,
zelikman2022star,
zhang2022automatic,
shao2023synthetic}.

In-context learning methods now routinely invoke tools, leading to LM pipelines that use retrieval models
\citep{
chen-etal-2017-reading,
lewis2020rag,
guu2020realm,
lazaridou2022internet,
izacard2022few},
multimodal foundation models, and more traditional tools like
APIs
\citep{nakano21webgpt} and
calculators. A number of toolkits %
have been developed to facilitate this, including
LangChain \citep{chase2022langchain},
Semantic Kernel \citep{semantickernel},
LlamaIndex \citep{llamaindex2022}, and many other retrieval and agent libraries.
These toolkits provide pre-packaged chains and agents that connect LMs with numerous accessible tools. However, they suffer from the pervasive prompt engineering challenges we address in DSPy: they express task-specific behavior through hand-written prompt templates (for detailed discussion, see Appendix~\ref{sec:dspy_vs_langchain}). %

Researchers are starting to apply discrete optimization and RL to find effective prompts, generally for a single logical LM call \citep{guo2023connecting,pryzant2023automatic,huang2022large,yang2023large}.
DSPy seeks to generalize this space: it offers a rich framework for optimizing \textit{arbitrary pipelines} from \textit{high-level declarative signatures}, by bootstrapping high-quality \textit{multi-stage demonstrations} with constraints. In this framework, DSPy teleprompters may apply optimization using model selection techniques like cross-validation or, in principle, with sophisticated techniques involving RL and LM feedback
\citep{hu2023enabling,zhao2023expel,shinn2023reflexion}
or learned or Bayesian hyperparameter optimization methods
\citep{bergstra2013making,akiba2019optuna}.

The present paper seeks to motivate DSPy as a programming model and to report new empirical findings from applying the DSPy compiler. This is inspired by formative work by
\citet{bergstra2010theano,bergstra2013making},
\citet{pytorch2019},
and
\citet{wolf2020transformers}, who
support their respective programming models with a mix of benchmark numbers and some qualitative measures. For the current paper, we focus on showing that DSPy and its compiler allow us to build outstanding LM systems without hand-crafted prompt strings, but instead from truly modular units, and that this opens up doors for systematically exploring a rich design space at a very high programmatic level of abstraction. %
